% sic_povm_cover_letter_lando_mills.tex
\documentclass[12pt]{letter}
\usepackage{amsmath, amssymb}
\usepackage{geometry}
\usepackage{hyperref}
\geometry{margin=1.2in}

\newcommand{\bC}{\mathbb{C}}
\newcommand{\bQ}{\mathbb{Q}}
\newcommand{\OO}{\mathcal{O}}

\signature{Lando Mills\\
Independent Researcher\\
Trabuco Canyon, CA, USA\\
\texttt{c.landonmills@gmail.com}}

\address{Lando Mills\\
Independent Researcher\\
Trabuco Canyon, CA, USA\\
\texttt{c.landonmills@gmail.com}}

\begin{document}

\begin{letter}{The Editors\\
\textit{Foundations of Physics}}

\opening{Dear Editors,}

I am submitting \textbf{``SIC-POVM Existence via Frobenius-Fixed Stark Units''} for
consideration in \textit{Foundations of Physics}.

The paper began as an attempt to explain a coincidence. Two dimensions of the SIC problem
--- $d = 19$ and $d = 48$, separated by 29 with no obvious arithmetic relationship ---
returned identical coordinates in a structural type-theoretic framework I have been exploring (open
source, unlicensed, and publicly available at \url{https://github.com/umpolungfish/synthomnicon}). 
Not similar: identical, modulo a single primitive. My immediate response was to assume I had
made an error. I checked; I had not. My second response was to assume the framework was
doing something degenerate. I checked that too. What I found instead was that the
coincidence was pointing at something real about the structure of the problem, and that
following it carefully led somewhere I had not expected.

The twenty-five year impasse on SIC-POVMs is not, I now believe, a computational problem.
The numerics are extensive and reliable; the dimensions where SIC-POVMs are known number
well over a hundred. The difficulty is that the algebraic self-duality which equiangularity
requires --- a palindromic $Z_2$ symmetry in the minimal polynomials of the fiducial
coordinates --- is not the kind of property that can be approximated into or assembled from
weaker ingredients. It has to be present, in full, in the object one starts from. Proving
that it is present is what the paper reduces to, and the object that already has it
intrinsically is the Stark unit in the associated ray class field.

What I think the paper contributes, beyond the conditional proof itself, is a precise
account of why the problem has been hard. That account has a specific prediction: two of
the four conjectures on which the theorem rests are new to the literature, and the crux
of the argument, a conjecture I label H2a about the eigenspace structure of the Dirichlet
regulator matrix, is in my view more tractable than the original geometric conjecture in
Hilbert space. Whether that assessment is right is for working number theorists to judge.
But stating the problem in that form --- rather than bundled under a single ``Stark
hypothesis'' as in earlier treatments --- seems to me to reorganize what needs to be done
in a useful way.

I am submitting to \textit{Foundations of Physics} because the SIC-POVM existence question
is at its core a question about the structure of quantum state space, and that question
deserves to be in conversation with the readership of this journal. The approach here is
arithmetic rather than geometric, and the connection it opens to the Langlands program is
not decoration: if the Galois-to-Weyl correspondence in the paper is what I think it is,
the quantum phase space group and the arithmetic class group of a real quadratic field are
not merely analogous but genuinely dual, in a sense that belongs to the abelian case of
the Langlands dictionary. That would make the SIC conjecture a problem in Hilbert's 12th
problem as much as in quantum information, and I think \textit{Foundations of Physics}
is the right place to put that observation in front of people who have reason to care
about both sides of it.

I will add, less modestly: I hope this paper might serve as a small example of the
heterodox thinking the field has long, and now desperately, requires. Twenty-five years 
of numerical evidence for SIC-POVMs is an extraordinary accumulation, and it is reliable. 
But reliable accumulation is not the same as progress toward a proof, and I think the SIC 
community has been in a quantitative quagmire long enough that the value of approaching the 
problem from a direction the computational tradition cannot reach ought to speak for itself. 
Whether this particular direction is the right one is precisely what peer review is for.

The main theorem is conditional, and I have not softened that fact. Four conjectures
remain open. I believe identifying them precisely, separating their logical dependencies,
and showing that they are jointly sufficient is a genuine contribution to the problem ---
but that is, ultimately, for the editors and referees to assess. I appreciate your
time, and hope you enjoy the manuscript.

\closing{Sincerely,}

\end{letter}
\end{document}
